% =============================================================================
% Draft mục 3.3.8 — Event Bus và Persistence
% =============================================================================

\subsection{Event Bus và Persistence}
\label{sec:event-bus-persistence}

Lớp persistence của hệ thống được thiết kế tách bạch hai mối quan tâm khác nhau. Trạng thái bền (durable state) phải sống sót qua restart tiến trình và là nguồn để rebuild in-memory cache; phần này được cài đặt bằng Postgres. Sự kiện thời gian thực (transient pub/sub) chỉ quan trọng khi đang truyền cho người dùng đang kết nối, không cần lưu trữ; phần này được cài đặt bằng một EventEmitter trong bộ nhớ.

Hai cơ chế này, đặt cạnh nhau, đóng vai trò tương đương với một event log Kafka trong các kiến trúc microservice phổ biến, nhưng với chi phí vận hành thấp hơn nhiều bậc — phù hợp với quy mô MVP và scope của luận văn. Phân tích so sánh được trình bày ở mục Phân tích thiết kế thay thế: Kafka vs Postgres + Emitter trong chương Đánh giá.

\subsubsection{Postgres và lược đồ ba bảng}

Cơ sở dữ liệu được truy cập qua Drizzle ORM (TypeScript-first, không codegen) với connection thuần từ \texttt{postgres-js}. Có ba bảng chính, dưới đây trình bày ngắn gọn (chi tiết cột và index xem \texttt{packages/backend/src/db/schema.ts}).

\subsubsection*{Bảng \texttt{intent\_orders}}

Bảng này có một dòng cho mỗi \texttt{IntentOrder}, với các cột chính gồm \texttt{order\_id} (primary key), \texttt{src\_chain}, \texttt{des\_chain}, \texttt{amount} (\texttt{text}, để giữ độ chính xác bigint), \texttt{incentive\_fee}, \texttt{deadline}, \texttt{created\_at}, \texttt{status} (chuỗi enum \texttt{QUEUED}/\texttt{PARTIAL}/\texttt{MATCHED}/\texttt{EXPIRED}) và \texttt{user\_address}. Bảng phục vụ rebuild orderbook khi \texttt{OrderStore.hydrate()} chạy lúc khởi động.

\subsubsection*{Bảng \texttt{match\_results}}

Đây là bảng append-only ghi mỗi sự kiện match, với các cột chính gồm \texttt{match\_id} (primary key), \texttt{matched\_at}, \texttt{orders} (JSONB --- mảng \texttt{MatchedOrderEntry}), \texttt{cycles} (JSONB) và \texttt{raw\_cycles} (JSONB --- snapshot cạnh từ thuật toán). Việc chọn JSONB cho ba trường mảng là vì hệ thống không bao giờ truy vấn vào trong các trường đó — chỉ ghi và đọc nguyên khối.

\subsubsection*{Bảng \texttt{executions}}

Bảng này lưu trạng thái thực thi mỗi \texttt{matchId}, với các cột \texttt{match\_id} (primary key), \texttt{status} (\texttt{pending}/\texttt{done}/\texttt{error}), \texttt{data} (JSONB) và \texttt{error}. Khác với hai bảng trên, dòng trong bảng này được cập nhật tại chỗ (qua \texttt{INSERT ... ON CONFLICT UPDATE}) khi Orchestrator chuyển trạng thái.

\subsubsection{Lazy initialization của DB client}

Vì module \texttt{db/client.ts} được \texttt{import} sớm (theo cơ chế hoisting của ES Module) trong khi \texttt{dotenv.config()} chỉ chạy sau khi tất cả import hoàn tất, một thiết kế đơn giản đọc \texttt{process.env.DATABASE\_URL} ở top-level sẽ throw lỗi.

Giải pháp được áp dụng là \emph{lazy initialization}. Module export một Proxy đại diện cho Drizzle handle; chỉ khi một thuộc tính của Proxy được truy cập lần đầu, hàm \texttt{getDb()} mới đọc biến môi trường, kết nối Postgres và cache instance lại. Cơ chế này làm cho \texttt{db/client.ts} không phụ thuộc thứ tự import, đồng thời cho phép unit test inject biến môi trường giả lập trước khi truy cập DB lần đầu.

\subsubsection{Mẫu write-through với fail-and-log}

Như đã mô tả ở \S\ref{sec:order-book}, mỗi mutation in-memory được mirror sang Postgres theo kiểu fire-and-forget. Cụ thể, Promise \texttt{insert/update} không được \texttt{await}, và handler \texttt{.catch(err => logDbError(scope, err))} đảm bảo lỗi không trở thành unhandled rejection và được log với namespace cụ thể (\texttt{orders.add}, \texttt{matchResults.add}, \texttt{executions.set}, ...).

Trade-off đã được phân tích ở \S\ref{sec:order-book}: nguy cơ lệch giữa in-memory và DB trong cửa sổ rất ngắn nếu tiến trình crash. Đây là chấp nhận được ở MVP, nhưng phương án thay thế thông qua \emph{outbox pattern} (append event vào một bảng riêng cùng transaction với mutation, sau đó dispatch tách biệt) được kiến nghị trong Future Work.

\subsubsection{\texttt{orderEvents}: in-memory event bus}

Ngoài durable state, hệ thống còn cần một kênh truyền sự kiện lifecycle thời gian thực đến các client đang kết nối qua Server-Sent Events (SSE). \texttt{orderEvents} là một \texttt{EventEmitter} của Node.js, sống trong cùng tiến trình backend và không persist.

Mỗi event là một đối tượng \texttt{OrderEvent} với cấu trúc gồm \texttt{orderId} (định danh order liên quan), \texttt{type} (một trong \texttt{queued | matched | plan | htlc\_created | withdrawn | done | error}), \texttt{message} (thông điệp ngắn cho UI), \texttt{data} (payload có cấu trúc, gồm txHash, chainKey, ...), và \texttt{timestamp} (Unix epoch).

Phía producer, \texttt{OrderStore} phát \texttt{queued} khi một order được \texttt{add()}, và phát \texttt{matched} khi \texttt{addMatchResult()}. \texttt{Orchestrator} phát \texttt{plan}, \texttt{htlc\_created}, \texttt{withdrawn}, \texttt{done}, \texttt{error} trong các bước thực thi của mình. Phía consumer, route \texttt{GET /api/intent/:orderId/events} mở một stream SSE: nó subscribe vào \texttt{orderEvents}, lọc theo \texttt{orderId} truyền vào, và push tuần tự xuống browser. Khi client đóng kết nối, listener được tháo gỡ. Với chỉ một consumer nội bộ (SSE handler), \texttt{EventEmitter} là đủ, không cần Redis pub/sub hay Kafka topic.

% [đã sửa: chuyển các \paragraph{...} và bullet list dạng "tên - mô tả" thành câu liền mạch trong văn xuôi]
